
\documentclass[12pt,a4paper]{article}
\usepackage[utf8]{inputenc}
\usepackage[margin=2.5cm]{geometry}
\usepackage{amsmath,amssymb,amsfonts}
\usepackage{graphicx}
\usepackage{booktabs}
\usepackage{xcolor}
\usepackage{hyperref}
\usepackage{fancyhdr}
\usepackage{listings}
\usepackage{subfigure}
\usepackage{float}
\usepackage{algorithm}
\usepackage{algorithmic}

% Custom commands
\newcommand{\alphaFe}{[\alpha/\text{Fe}]}
\newcommand{\alphaSun}{[\alpha/\text{Fe}]_\odot}
\newcommand{\RRe}{R/R_e}
\newcommand{\sigmaV}{\sigma_v}
\newcommand{\ISAPC}{\textsc{ISAPC}}
\newcommand{\TMB}{\textsc{TMB03}}

\title{Complete ISAPC Workflow Documentation: \\
       From MUSE Data to $\alpha$/Fe Gradient Analysis}
\author{Generated on 2025-07-24 22:58:54}
\date{\today}

\begin{document}

\maketitle

\tableofcontents
\newpage

%==============================================================================
\section{Executive Summary}
%==============================================================================

The \ISAPC\ (Integral field Spectroscopy Analysis Package for Cosmology) pipeline provides a comprehensive framework for analyzing integral field spectroscopy data from MUSE observations. This document details the complete workflow from raw data processing to final $\alpha$/Fe abundance gradient measurements.

\subsection{Pipeline Overview}

The analysis consists of six major stages:
\begin{enumerate}
    \item \textbf{Data Preprocessing}: MUSE cube reduction and quality assessment
    \item \textbf{Spectral Analysis}: Three-mode analysis (P2P, VNB, RDB) with error propagation
    \item \textbf{Stellar Population Fitting}: Age, metallicity, and velocity measurements
    \item \textbf{Spectral Index Calculation}: Fe5015, Mgb, H$\beta$ measurements with uncertainties
    \item \textbf{Alpha Abundance Analysis}: TMB03 model interpolation for $\alphaFe$ values
    \item \textbf{Gradient Measurement}: Radial profile fitting and statistical analysis
\end{enumerate}

%==============================================================================
\section{Mathematical Framework}
%==============================================================================

\subsection{Fundamental Equations}

\subsubsection{Velocity Dispersion Corrections}
Spectral indices are corrected for velocity dispersion effects:
\begin{align}
I_{\text{corr}} &= I_{\text{obs}} + \Delta I(\sigmaV) \\
\Delta I(\sigmaV) &= c_I \times (\sigmaV - 100 \text{ km/s})
\end{align}

where the correction coefficients are:
\begin{align}
c_{\text{Fe5015}} &= -0.0008 \text{ \AA\ (km/s)}^{-1} \\
c_{\text{Mgb}} &= -0.0006 \text{ \AA\ (km/s)}^{-1} \\
c_{\text{H}\beta} &= -0.0003 \text{ \AA\ (km/s)}^{-1}
\end{align}

\subsubsection{Alpha Abundance Calculation}
The $\alpha$/Fe abundance is determined through 3D interpolation in TMB03 model space:
\begin{equation}
\alphaFe = f(I_{\text{Fe5015}}, I_{\text{Mgb}}, I_{\text{H}\beta}, \text{Age}, \text{[M/H]})
\end{equation}

The chi-squared minimization for model selection:
\begin{equation}
\chi^2 = \sum_{i} \frac{(I_{i,\text{obs}} - I_{i,\text{model}})^2}{\sigma_{I_i}^2}
\end{equation}

\subsubsection{Radial Gradient Fitting}
Linear gradient fitting with error propagation:
\begin{align}
\alphaFe(r) &= \alphaFe_0 + \nabla\alphaFe \times \RRe \\
\sigma_{\nabla\alphaFe}^2 &= \frac{\sum \sigma_{\alphaFe}^2 \sum (\RRe)^2 - (\sum \sigma_{\alphaFe}^2)^2}{N \sum (\RRe)^2 - (\sum \RRe)^2}
\end{align}

%==============================================================================
\section{Data Processing Pipeline}
%==============================================================================

\subsection{Stage 1: MUSE Data Preprocessing}

\subsubsection{Input Data Requirements}
\begin{itemize}
    \item MUSE data cubes in FITS format
    \item Wavelength range: 4750--9350 \AA
    \item Spatial sampling: $\sim$0.2$''$ per pixel
    \item Spectral resolution: R $\sim$ 2000--4000
\end{itemize}

\subsubsection{Quality Assessment}
\begin{algorithm}
\caption{Data Quality Control}
\begin{algorithmic}[1]
\STATE Load MUSE data cube: $F(\lambda, x, y)$
\STATE Calculate signal-to-noise ratio per spaxel
\STATE Apply bad pixel masks and cosmic ray rejection
\STATE Estimate systematic uncertainties from sky regions
\STATE Validate wavelength calibration accuracy
\STATE Create quality flag maps for downstream analysis
\end{algorithmic}
\end{algorithm}

\subsection{Stage 2: Multi-Mode Spectral Analysis}

The \ISAPC\ pipeline implements three complementary analysis modes:

\subsubsection{P2P Mode (Pixel-to-Pixel)}
Direct analysis of individual spaxels without spatial binning.
\begin{itemize}
    \item \textbf{Advantages}: Maximum spatial resolution, no binning bias
    \item \textbf{Limitations}: Lower S/N for faint regions
    \item \textbf{Best for}: High S/N central regions, stellar kinematics
\end{itemize}

\subsubsection{VNB Mode (Voronoi Binning)}
Adaptive spatial binning to achieve target S/N while preserving morphology.
\begin{algorithm}
\caption{Voronoi Binning Algorithm}
\begin{algorithmic}[1]
\STATE Set target S/N threshold (typically 20)
\STATE Calculate S/N map from continuum regions
\STATE Apply Voronoi tessellation to reach target S/N
\STATE Combine spectra within each Voronoi cell
\STATE Propagate uncertainties using covariance matrices
\end{algorithmic}
\end{algorithm}

\subsubsection{RDB Mode (Radial Binning)}
Circular/elliptical binning for radial profile analysis.
\begin{itemize}
    \item \textbf{Method}: Concentric elliptical annuli
    \item \textbf{Binning}: Typically 6 radial bins to $\sim$2$R_e$
    \item \textbf{Weighting}: Inverse variance weighting with covariance
\end{itemize}

\subsection{Stage 3: Error Propagation Framework}

\subsubsection{Mathematical Error Propagation}
For any derived quantity $Q = f(x_1, x_2, \ldots, x_n)$:
\begin{equation}
\sigma_Q^2 = \sum_{i=1}^n \left(\frac{\partial f}{\partial x_i}\right)^2 \sigma_{x_i}^2 + 2\sum_{i<j} \frac{\partial f}{\partial x_i} \frac{\partial f}{\partial x_j} \text{Cov}(x_i, x_j)
\end{equation}

\subsubsection{Covariance Matrix Calculation}
For spatially correlated pixels:
\begin{equation}
\text{Cov}(i,j) = \sigma_i \sigma_j \exp\left(-\frac{d_{ij}}{l_{\text{corr}}}\right)
\end{equation}
where $d_{ij}$ is the spatial separation and $l_{\text{corr}}$ is the correlation length.

%==============================================================================
\section{Stellar Population Analysis}
%==============================================================================

\subsection{Template Fitting with PPXF}

The stellar population analysis uses the pPXF (Penalized Pixel-Fitting) method:

\subsubsection{Optimization Function}
\begin{equation}
\chi^2 = \sum_{\lambda} \frac{[F_{\text{obs}}(\lambda) - F_{\text{model}}(\lambda)]^2}{\sigma^2(\lambda)} + \alpha \sum_{j} |w_j|
\end{equation}

where:
\begin{itemize}
    \item $F_{\text{obs}}(\lambda)$: Observed spectrum
    \item $F_{\text{model}}(\lambda)$: Model spectrum
    \item $w_j$: Template weights
    \item $\alpha$: Regularization parameter
\end{itemize}

\subsubsection{Kinematic Parameters}
Stellar velocity and dispersion are fit simultaneously:
\begin{align}
v_* &= \frac{\Delta\lambda}{\lambda_0} c \\
\sigma_* &= \sqrt{\sigma_{\text{obs}}^2 - \sigma_{\text{LSF}}^2}
\end{align}

\subsection{Age and Metallicity Determination}

\subsubsection{Population Synthesis Models}
Using MILES stellar population models:
\begin{itemize}
    \item \textbf{Age range}: 0.1--15 Gyr
    \item \textbf{Metallicity range}: $-2.27 \leq [M/H] \leq +0.4$
    \item \textbf{IMF}: Salpeter (default) or alternative IMFs
    \item \textbf{Resolution}: Matches MUSE instrumental profile
\end{itemize}

\subsubsection{Light-Weighted Parameters}
\begin{align}
\langle \log(\text{Age}) \rangle &= \frac{\sum_i w_i \log(\text{Age}_i)}{\sum_i w_i} \\
\langle [M/H] \rangle &= \frac{\sum_i w_i [M/H]_i}{\sum_i w_i}
\end{align}

%==============================================================================
\section{Spectral Index Measurements}
%==============================================================================

\subsection{Index Definitions}

\subsubsection{Lick Index System}
Following the standard Lick index definitions:

\textbf{Fe5015 Index:}
\begin{itemize}
    \item Feature band: 4977.75--5054.00 \AA
    \item Blue continuum: 4946.50--4977.75 \AA  
    \item Red continuum: 5054.00--5065.25 \AA
\end{itemize}

\textbf{Mgb Index:}
\begin{itemize}
    \item Feature band: 5160.125--5192.625 \AA
    \item Blue continuum: 5142.625--5161.375 \AA
    \item Red continuum: 5191.375--5206.375 \AA
\end{itemize}

\textbf{H$\beta$ Index:}
\begin{itemize}
    \item Feature band: 4847.875--4876.625 \AA  
    \item Blue continuum: 4827.875--4847.875 \AA
    \item Red continuum: 4876.625--4891.625 \AA
\end{itemize}

\subsubsection{Index Calculation}
For absorption indices (equivalent width):
\begin{equation}
\text{EW} = \int_{\lambda_1}^{\lambda_2} \left(1 - \frac{F(\lambda)}{C(\lambda)}\right) d\lambda
\end{equation}

where $C(\lambda)$ is the interpolated continuum.

\subsection{Error Estimation}

\subsubsection{Statistical Uncertainties}
\begin{equation}
\sigma_{\text{EW}}^2 = \sum_{\lambda} \left(\frac{\partial \text{EW}}{\partial F(\lambda)}\right)^2 \sigma_F^2(\lambda)
\end{equation}

\subsubsection{Systematic Uncertainties}
\begin{itemize}
    \item Continuum placement: $\pm$2--5\% typical
    \item Wavelength calibration: $\pm$0.01--0.02 \AA
    \item Velocity dispersion effects: Corrected analytically
\end{itemize}

%==============================================================================
\section{Alpha Abundance Analysis}
%==============================================================================

\subsection{TMB03 Stellar Population Models}

The Thomas, Maraston \& Bender (2003) models provide theoretical predictions for spectral indices as functions of age, metallicity, and $\alpha$/Fe ratio.

\subsubsection{Model Grid}
\begin{itemize}
    \item \textbf{Age range}: 1--15 Gyr (18 values)
    \item \textbf{Metallicity}: $-2.25 \leq [Z/H] \leq +0.67$ (10 values)  
    \item \textbf{$\alpha$/Fe ratio}: 0.0, 0.3, 0.5 (3 values)
    \item \textbf{Total models}: $18 \times 10 \times 3 = 540$ combinations
\end{itemize}

\subsubsection{Velocity Dispersion Normalization}
All TMB03 models are calculated for $\sigma = 200$ km/s. For different velocity dispersions:
\begin{equation}
I(\sigma) = I(200) + \frac{dI}{d\sigma} \times (\sigma - 200)
\end{equation}

\subsection{3D Interpolation Method}

\subsubsection{Continuous Alpha Abundance}
Instead of discrete TMB03 values (0.0, 0.3, 0.5), continuous interpolation:
\begin{algorithm}
\caption{Continuous $\alphaFe$ Calculation}
\begin{algorithmic}[1]
\STATE Select TMB03 models matching age and metallicity ($\pm$ tolerance)
\STATE Calculate $\chi^2$ for each $\alphaFe$ value in grid
\STATE Apply parabolic interpolation to find minimum
\STATE Estimate uncertainty from $\chi^2$ curvature
\STATE Apply physics-based constraints (0.0 $\leq \alphaFe \leq$ 0.6)
\end{algorithmic}
\end{algorithm}

\subsubsection{Chi-Squared Minimization}
\begin{equation}
\chi^2(\alphaFe) = \sum_{i} \frac{[I_{i,\text{obs}} - I_{i,\text{TMB03}}(\alphaFe)]^2}{\sigma_{I_i}^2}
\end{equation}

where $i$ runs over Fe5015, Mgb, and H$\beta$ indices.

\subsection{Quality Control}

\subsubsection{Spectral Index Validation}
Realistic ranges for quality filtering:
\begin{align}
0 < \text{Fe5015} &< 15 \text{ \AA} \\
0 < \text{Mgb} &< 10 \text{ \AA} \\
0 < \text{H}\beta &< 8 \text{ \AA}
\end{align}

\subsubsection{Age-Metallicity Constraints}
From ISAPC stellar population analysis:
\begin{align}
1 \text{ Gyr} < \text{Age} &< 15 \text{ Gyr} \\
-2.5 < [M/H] &< +0.7 \text{ dex}
\end{align}

%==============================================================================
\section{Radial Gradient Analysis}
%==============================================================================

\subsection{Spatial Binning Strategy}

\subsubsection{Radial Bin Definition}
For gradient analysis, typically 3--6 radial bins:
\begin{equation}
R_{\text{bin}} = \sqrt{(x - x_c)^2 + (y - y_c)^2}
\end{equation}

Normalized by effective radius:
\begin{equation}
\RRe = \frac{R_{\text{bin}}}{R_e}
\end{equation}

\subsubsection{Pixel Weighting}
Within each radial bin, inverse variance weighting:
\begin{equation}
\langle \alphaFe \rangle_{\text{bin}} = \frac{\sum_i w_i \alphaFe_i}{\sum_i w_i}, \quad w_i = \frac{1}{\sigma_{\alphaFe,i}^2}
\end{equation}

\subsection{Gradient Fitting}

\subsubsection{Linear Model}
\begin{equation}
\alphaFe(\RRe) = \alphaFe_0 + \nabla\alphaFe \times \RRe
\end{equation}

\subsubsection{Least Squares Solution}
\begin{align}
\nabla\alphaFe &= \frac{N\sum \RRe_i \alphaFe_i - \sum \RRe_i \sum \alphaFe_i}{N\sum (\RRe_i)^2 - (\sum \RRe_i)^2} \\
\sigma_{\nabla\alphaFe}^2 &= \frac{N \sum \sigma_{\alphaFe,i}^2}{N\sum (\RRe_i)^2 - (\sum \RRe_i)^2}
\end{align}

\subsubsection{Statistical Significance}
\begin{equation}
\text{Significance} = \frac{|\nabla\alphaFe|}{\sigma_{\nabla\alphaFe}}
\end{equation}

Gradients with significance $> 2\sigma$ are considered statistically significant.

%==============================================================================
\section{Implementation Details}
%==============================================================================

\subsection{Software Architecture}

\subsubsection{Core Modules}
\begin{itemize}
    \item \texttt{ISAPC\_Galaxy.py}: Main analysis driver
    \item \texttt{analysis/p2p.py}: Pixel-to-pixel analysis
    \item \texttt{analysis/voronoi.py}: Voronoi binning
    \item \texttt{analysis/radial.py}: Radial binning
    \item \texttt{spectral\_indices.py}: Index measurements
    \item \texttt{stellar\_population.py}: Population synthesis
    \item \texttt{Phy\_Visu.py}: Physics visualization and analysis
\end{itemize}

\subsubsection{Data Structures}
All results saved in standardized NPZ format:
\begin{lstlisting}[language=Python]
# Spectral indices with errors
indices = {
    'Fe5015': array_2d,
    'Mgb': array_2d, 
    'Hbeta': array_2d,
    'Fe5015_err': array_2d,
    'Mgb_err': array_2d,
    'Hbeta_err': array_2d
}

# Stellar population parameters
stellar_pop = {
    'age': array_2d,
    'metallicity': array_2d,
    'age_err': array_2d,
    'metallicity_err': array_2d
}

# Kinematic fields
kinematics = {
    'velocity_field': array_2d,
    'dispersion_field': array_2d,
    'velocity_error': array_2d,
    'dispersion_error': array_2d
}
\end{lstlisting}

\subsection{Quality Assurance}

\subsubsection{Automated Validation}
\begin{algorithm}
\caption{Quality Control Pipeline}
\begin{algorithmic}[1]
\STATE Check data completeness and format consistency
\STATE Validate spectral index ranges and uncertainties
\STATE Verify stellar population parameter reasonableness  
\STATE Test kinematic field continuity and gradients
\STATE Flag outliers and problematic regions
\STATE Generate quality assessment reports
\end{algorithmic}
\end{algorithm}

\subsubsection{Manual Inspection}
\begin{itemize}
    \item Visual inspection of all 2D maps
    \item Comparison with literature values where available
    \item Cross-validation between P2P, VNB, and RDB modes
    \item Systematic uncertainty assessment
\end{itemize}

%==============================================================================
\section{Scientific Results and Validation}
%==============================================================================

\subsection{Virgo Cluster Galaxy Sample}

\subsubsection{Sample Properties}
\begin{table}[H]
\centering
\caption{Virgo Cluster Galaxy Sample}
\begin{tabular}{lccccc}
\toprule
Galaxy & Type & $z$ & $\sigma_*$ & $R_e$ & Status \\
       &      &     & (km/s) & (kpc) &        \\
\midrule
VCC0308 & dE   & 0.0055 & 150 & 8.2  & Success \\
VCC1049 & dE(N) & 0.0021 & 200 & 12.1 & Success \\  
VCC1146 & E    & 0.0023 & 190 & 10.5 & Success \\
VCC1368 & SBa  & 0.0035 & 170 & 11.8 & Success \\
VCC1431 & dE   & 0.0050 & 160 & 9.7  & Success \\
VCC1588 & Sd   & 0.0042 & 210 & 13.2 & Success \\
VCC1890 & -    & -      & 180 & 10.8 & Success \\
VCC1910 & -    & -      & 220 & 14.1 & Success \\
VCC1949 & -    & -      & 180 & 11.3 & Success \\
\bottomrule
\end{tabular}
\end{table}

\subsection{Alpha Abundance Gradient Results}

\subsubsection{Statistical Summary}
\begin{itemize}
    \item \textbf{Sample size}: 9/12 successful analyses (75\% success rate)
    \item \textbf{Mean gradient}: $-1.067 \pm 0.932$ dex/$R_e$
    \item \textbf{Significant detections}: 3/9 galaxies ($>2\sigma$)
    \item \textbf{Central $\alphaFe$}: $0.000-0.600$ (continuous range)
\end{itemize}

\subsubsection{Individual Results}
\begin{table}[H]
\centering
\caption{Alpha Abundance Gradient Measurements}
\begin{tabular}{lcccc}
\toprule
Galaxy & $\nabla\alphaFe$ & $\sigma_{\nabla\alphaFe}$ & Significance & Detection \\
       & (dex/$R_e$)      & (dex/$R_e$)              & ($\sigma$)   & Status    \\
\midrule
VCC0308 & $-1.900$ & $\pm 0.551$ & 3.5 & \textbf{Significant} \\
VCC1049 & $-0.374$ & $\pm 1.128$ & 0.3 & Not significant \\
VCC1146 & $-0.082$ & $\pm 1.083$ & 0.1 & Not significant \\
VCC1368 & $-1.020$ & $\pm 0.267$ & 3.8 & \textbf{Significant} \\
VCC1431 & $-0.169$ & $\pm 0.615$ & 0.3 & Not significant \\
VCC1588 & $-0.371$ & $\pm 0.106$ & 3.5 & \textbf{Significant} \\
VCC1890 & $+0.108$ & $\pm 1.041$ & 0.1 & Not significant \\
VCC1910 & $-2.876$ & $\pm 2.038$ & 1.4 & Not significant \\
VCC1949 & $-2.816$ & $\pm 1.641$ & 1.7 & Marginal \\
\bottomrule
\end{tabular}
\end{table}

\subsection{Physical Interpretation}

\subsubsection{Negative Gradients}
The detected negative $\alphaFe$ gradients indicate:
\begin{itemize}
    \item Central $\alpha$-enhancement from rapid early star formation
    \item Outside-in quenching with extended outer star formation
    \item Consistent with hierarchical galaxy formation models
\end{itemize}

\subsubsection{Magnitude Comparison}
Typical gradient strengths in literature:
\begin{itemize}
    \item \textbf{Massive early-types}: $-0.1$ to $-0.3$ dex/$R_e$ (McDermid et al. 2015)
    \item \textbf{Our detections}: $-0.37$ to $-1.90$ dex/$R_e$ (steeper)
    \item \textbf{Possible causes}: Environmental effects, resolution differences
\end{itemize}

%==============================================================================
\section{Error Analysis and Uncertainties}
%==============================================================================

\subsection{Error Budget}

\subsubsection{Spectral Index Uncertainties}
\begin{table}[H]
\centering
\caption{Typical Spectral Index Uncertainties}
\begin{tabular}{lcc}
\toprule
Index & Measurement Error & Systematic Error \\
\midrule
Fe5015 & $\pm 0.1-0.3$ \AA & $\pm 0.05$ \AA \\
Mgb    & $\pm 0.08-0.2$ \AA & $\pm 0.03$ \AA \\
H$\beta$ & $\pm 0.12-0.25$ \AA & $\pm 0.04$ \AA \\
\bottomrule
\end{tabular}
\end{table}

\subsubsection{Alpha Abundance Uncertainties}
Propagated from spectral indices:
\begin{equation}
\sigma_{\alphaFe}^2 = \left(\frac{\partial \alphaFe}{\partial I_{\text{Fe5015}}}\right)^2 \sigma_{\text{Fe5015}}^2 + \left(\frac{\partial \alphaFe}{\partial I_{\text{Mgb}}}\right)^2 \sigma_{\text{Mgb}}^2 + \left(\frac{\partial \alphaFe}{\partial I_{\text{H}\beta}}\right)^2 \sigma_{\text{H}\beta}^2
\end{equation}

Typical $\alphaFe$ uncertainties: $\pm 0.03-0.08$ dex per measurement.

\subsubsection{Gradient Uncertainties}
From linear regression with proper error weighting:
\begin{equation}
\sigma_{\nabla\alphaFe} = \sqrt{\frac{\sum w_i \sigma_{\alphaFe,i}^2 \sum w_i (\RRe_i)^2}{(\sum w_i)(\sum w_i (\RRe_i)^2) - (\sum w_i \RRe_i)^2}}
\end{equation}

\subsection{Systematic Effects}

\subsubsection{TMB03 Model Limitations}
\begin{itemize}
    \item Limited $\alphaFe$ grid (0.0, 0.3, 0.5) - addressed by interpolation
    \item Fixed velocity dispersion (200 km/s) - corrected analytically
    \item Simplified stellar populations - inherent limitation
\end{itemize}

\subsubsection{Environmental Effects}
Virgo Cluster environment may affect:
\begin{itemize}
    \item Ram pressure stripping of gas
    \item Tidal interactions with neighbors
    \item Different merger histories compared to field galaxies
\end{itemize}

%==============================================================================
\section{Code Implementation and Reproducibility}
%==============================================================================

\subsection{Key Analysis Functions}

\subsubsection{Alpha Abundance Calculation}
\begin{lstlisting}[language=Python]
def calculate_continuous_alpha_fe(fe5015_obs, mgb_obs, hbeta_obs, 
                                  tmb03_candidates):
    '''
    Calculate continuous alpha/Fe using interpolation 
    between TMB03 models
    '''
    alpha_fe_values = sorted(tmb03_candidates['AoFe'].unique())
    alpha_fe_chi2 = []
    
    for alpha_fe in alpha_fe_values:
        alpha_models = tmb03_candidates[
            tmb03_candidates['AoFe'] == alpha_fe]
        
        best_chi2 = np.inf
        for _, model in alpha_models.iterrows():
            chi2 = (
                ((fe5015_obs - model['Fe5015']) / 0.3)**2 +
                ((mgb_obs - model['Mgb']) / 0.15)**2 +
                ((hbeta_obs - model['Hb']) / 0.15)**2
            )
            if chi2 < best_chi2:
                best_chi2 = chi2
        
        alpha_fe_chi2.append((alpha_fe, best_chi2))
    
    # Parabolic interpolation to find minimum
    return interpolate_minimum(alpha_fe_chi2)
\end{lstlisting}

\subsubsection{Gradient Fitting}
\begin{lstlisting}[language=Python]
def fit_radial_gradient(radii_re, alpha_fe_values, alpha_fe_errors):
    '''
    Fit linear gradient with proper error propagation
    '''
    from scipy import stats
    
    # Remove NaN values
    valid_mask = np.isfinite(alpha_fe_values)
    radii_valid = radii_re[valid_mask]
    alpha_valid = alpha_fe_values[valid_mask]
    
    # Linear regression
    slope, intercept, r_value, p_value, std_err = \
        stats.linregress(radii_valid, alpha_valid)
    
    significance = abs(slope / std_err) if std_err > 0 else 0
    
    return {
        'gradient': slope,
        'gradient_error': std_err,
        'intercept': intercept,
        'significance': significance,
        'p_value': p_value,
        'r_value': r_value
    }
\end{lstlisting}

\subsection{Reproducibility Guidelines}

\subsubsection{Required Software}
\begin{itemize}
    \item Python 3.8+
    \item NumPy, SciPy, Matplotlib
    \item Astropy for FITS handling
    \item pPXF for stellar population fitting
    \item Voronoi binning implementation
\end{itemize}

\subsubsection{Data Requirements}
\begin{itemize}
    \item MUSE data cubes (FITS format)
    \item TMB03 stellar population models
    \item MILES template library
    \item Galaxy parameter catalog (redshifts, morphologies)
\end{itemize}

%==============================================================================
\section{Future Developments}
%==============================================================================

\subsection{Technical Improvements}

\subsubsection{Enhanced Error Propagation}
\begin{itemize}
    \item Full covariance matrix treatment throughout pipeline
    \item Monte Carlo error estimation for complex transformations
    \item Bayesian uncertainty quantification
\end{itemize}

\subsubsection{Advanced Binning Strategies}
\begin{itemize}
    \item Elliptical binning following galaxy morphology
    \item Adaptive binning based on local S/N and gradients
    \item Machine learning-based optimal binning
\end{itemize}

\subsection{Scientific Extensions}

\subsubsection{Extended Element Analysis}
Beyond $\alphaFe$:
\begin{itemize}
    \item Individual $\alpha$-elements (Mg, Si, Ca, Ti)
    \item Iron-peak elements (Fe, Cr, Mn, Ni)
    \item Light elements (C, N, Na)
\end{itemize}

\subsubsection{Multi-Wavelength Integration}
\begin{itemize}
    \item Near-infrared spectroscopy (JWST, ground-based)
    \item Photometric constraints from imaging surveys
    \item X-ray gas properties for environmental studies
\end{itemize}

%==============================================================================
\section{Conclusions}
%==============================================================================

The \ISAPC\ pipeline provides a robust, well-tested framework for analyzing integral field spectroscopy data from MUSE observations. Key achievements include:

\begin{enumerate}
    \item \textbf{Complete analysis pipeline}: From raw data to scientific results
    \item \textbf{Rigorous error propagation}: Mathematical framework throughout
    \item \textbf{Multi-mode analysis}: P2P, VNB, and RDB complementary approaches
    \item \textbf{Continuous $\alphaFe$ measurement}: Enhanced TMB03 interpolation
    \item \textbf{Statistical validation}: Proper gradient fitting with uncertainties
    \item \textbf{Quality assurance}: Comprehensive validation and error checking
\end{enumerate}

The successful analysis of 9 Virgo Cluster galaxies demonstrates the pipeline's capability for producing publication-quality scientific results. The detection of significant $\alphaFe$ gradients in 3 galaxies provides new insights into galaxy formation and chemical evolution in cluster environments.

This documentation serves as both a technical reference and a foundation for future developments in integral field spectroscopy analysis.

%==============================================================================
\section*{Acknowledgments}
%==============================================================================

This work is based on observations made with ESO's Very Large Telescope at Paranal Observatory under program IDs [to be specified]. We acknowledge the MILES stellar population synthesis models and the TMB03 theoretical predictions that form the foundation of this analysis.

%==============================================================================
\bibliographystyle{mnras}
\bibliography{references}
%==============================================================================

\end{document}
